\documentclass[a4paper,11pt]{ctexart}
\usepackage{amsmath, amssymb, amsfonts}
\usepackage{geometry}
\geometry{left=1.5cm, right=1.5cm, top=2.0cm, bottom=2.0cm, headsep=0.3cm, footskip=1cm}
\usepackage{listings}
\usepackage{xcolor}
\usepackage{tikz}
\usepackage{pgfplots}
\usepackage{hyperref}
\usepackage{fancyhdr}
\usepackage{tcolorbox}
\tcbuselibrary{breakable, skins}
\usepackage{array}
\usepackage{booktabs}
\usepackage[shortlabels]{enumitem}
\usepackage{needspace}
\usepackage{multirow}

\AtBeginDocument{
  \hypersetup{colorlinks=true, linkcolor=blue, filecolor=magenta, urlcolor=cyan}
}

\setlist{nosep, itemsep=0.8ex, parsep=0.1em}
\setlist[itemize]{leftmargin=1.8em}
\setlist[enumerate]{leftmargin=1.8em}

% ===== 颜色定义 =====
\definecolor{codegreen}{rgb}{0,0.6,0}
\definecolor{codegray}{rgb}{0.5,0.5,0.5}
\definecolor{codepurple}{rgb}{0.58,0,0.82}
\definecolor{codebg}{rgb}{0.98,0.98,0.98}
\definecolor{tipsblue}{rgb}{0.85,0.93,1.0}
\definecolor{highlightyellow}{rgb}{1.0,0.97,0.8}

% ===== 代码块样式 =====
\lstdefinestyle{mystyle}{
    backgroundcolor=\color{codebg},
    commentstyle=\color{codegreen},
    keywordstyle=\color{magenta}\bfseries,
    numberstyle=\tiny\color{codegray},
    stringstyle=\color{codepurple},
    basicstyle=\ttfamily\small,
    breakatwhitespace=false,
    breaklines=true,
    captionpos=b,
    keepspaces=true,
    numbers=left,
    numbersep=6pt,
    showspaces=false,
    showstringspaces=false,
    showtabs=false,
    tabsize=2,
    language=Python,
    frame=single,
    framerule=0.5pt,
    rulecolor=\color{gray!40}
}
\lstset{style=mystyle}

% ===== 彩色框 =====
\newtcolorbox{problembox}[1][]{
    colback=white,
    colframe=blue!60!black,
    title=\textbf{#1},
    fonttitle=\bfseries\small,
    boxrule=1pt,
    arc=3pt, outer arc=3pt,
    coltitle=black,
    before skip=10pt, after skip=8pt
}

\newtcolorbox{solutionbox}[1][]{
    enhanced, breakable,
    colback=green!3!white,
    colframe=green!50!black,
    title=\textbf{#1},
    fonttitle=\bfseries\small,
    boxrule=1pt,
    arc=3pt, outer arc=3pt,
    coltitle=black,
    before skip=8pt, after skip=10pt
}

\newtcolorbox{metaphorbox}[1][]{
    colback=orange!5!white,
    colframe=orange!70!black,
    title=\textbf{#1},
    fonttitle=\bfseries\small,
    boxrule=1pt,
    arc=3pt, outer arc=3pt,
    coltitle=black,
    before skip=8pt, after skip=8pt
}

\newtcolorbox{beginnerbox}[1][]{
    colback=tipsblue,
    colframe=blue!50!black,
    title=\textbf{#1},
    fonttitle=\bfseries\small,
    boxrule=1pt,
    arc=3pt, outer arc=3pt,
    coltitle=black,
    before skip=8pt, after skip=8pt
}

\newtcolorbox{appbox}[1][]{
    colback=green!3!white,
    colframe=green!60!black,
    title=\textbf{#1},
    fonttitle=\bfseries\small,
    boxrule=1pt,
    arc=3pt, outer arc=3pt,
    coltitle=black,
    before skip=8pt, after skip=10pt
}

\newtcolorbox{keybox}[1][]{
    colback=highlightyellow,
    colframe=orange!80!black,
    title=\textbf{#1},
    fonttitle=\bfseries\small,
    boxrule=1pt,
    arc=3pt, outer arc=3pt,
    coltitle=black,
    before skip=8pt, after skip=8pt
}

% ===== 页眉页脚 =====
\pagestyle{fancy}
\fancyhf{}
\fancyhead[R]{深度学习-第4章}
\fancyhead[L]{反向传播算法 · 练习题}
\fancyfoot[C]{\thepage}
\addtolength{\headheight}{2pt}

\parskip=2pt plus 1pt minus 0.5pt
\linespread{1.35}
\raggedbottom
\widowpenalty=10000
\clubpenalty=10000

\title{\textbf{深度学习\\第4章\quad 反向传播算法}\\[0.5em]
{\large\color{gray}——神经网络是怎么"知道"自己哪里错了，又是怎么改正的？}}
\date{\today}

\begin{document}

\maketitle

% ===== 章节导读 =====
\begin{beginnerbox}[本章题目导览：你将挑战哪些核心问题？]
    反向传播是深度学习的心脏——没有它，网络就无法从错误中学习。本章 6 道题从计算图的"画法"一路打通到完整代码实现，把这个被误解最多的算法彻底讲清楚：

    \begin{itemize}
        \item \textbf{Q4-01（计算图阅读）}：前向传播的计算图构建——给定一个数学表达式，画出对应的计算图，标注中间变量，为反向传播打基础。
        \item \textbf{Q4-02（链式法则推导）}：链式法则手动推导——从复合函数的角度，一步步写出梯度沿计算图反向流动的公式，把链式法则从"记公式"变成"会计算"。
        \item \textbf{Q4-03（节点反向传播）}：加法节点与乘法节点的反向传播——两类最基本的节点，梯度各自怎么传？填表 + 数值验证。
        \item \textbf{Q4-04（激活层推导）}：ReLU 和 Sigmoid 的反向传播——推导两种激活函数的梯度公式，说明为什么 ReLU 的梯度"更干净"。
        \item \textbf{Q4-05（Affine 层推导）}：矩阵形式的反向传播——全连接层（Affine 层）在批量输入下，权重矩阵 $W$、偏置 $b$、输入 $X$ 的梯度各是什么形状？如何推导？
        \item \textbf{Q4-06（综合代码实现）}：完整反向传播代码追踪——阅读一个含 Affine + ReLU + Softmax-CrossEntropy 的两层网络实现，追踪梯度在各层之间的流动。
    \end{itemize}
\end{beginnerbox}

\section{第 4 章\quad 反向传播算法}

上一章我们理解了神经网络的前向传播——数据从输入层经过若干变换流向输出层，得到预测值，和标签比较，得到一个 Loss。

但"知道 Loss 有多大"只完成了一半。真正的学习发生在之后：

\begin{itemize}
    \item \textbf{Loss 对哪个权重的影响最大？}——这些权重需要调整最多；
    \item \textbf{Loss 对哪个权重的影响接近零？}——这些权重暂时不用动；
    \item \textbf{应该让权重变大还是变小？}——朝着让 Loss 减小的方向走。
\end{itemize}

回答这三个问题，需要计算 Loss 对每个参数的偏导数（梯度）。对于一个含数百层、数百万参数的网络，手推梯度几乎不可能——\textbf{反向传播算法}就是高效自动完成这件事的核心工具。

\begin{metaphorbox}[先建立一个总感觉：反向传播是什么？]
    \textbf{把神经网络比作一条装配流水线，Loss 是最终的"质检报告"。}

    \begin{itemize}
        \item \textbf{前向传播} = 零件从左到右经过每道工序，最终组装出产品；
        \item \textbf{Loss} = 质检员发现产品有缺陷，给出一个"不合格分数"；
        \item \textbf{反向传播} = 质检报告从右向左传递：先告诉最后一道工序"你的错误占了多少比例"，再告诉倒数第二道工序，……一直追溯到源头；
        \item \textbf{梯度下降} = 每道工序根据自己被"归责"的比例，调整操作参数，下次少出错。
    \end{itemize}

    反向传播的本质就是：\textbf{用链式法则，把"最终错误"按照因果比例，分配给每一个参数。}
\end{metaphorbox}


%% ============================================================
\Needspace{5\baselineskip}
\begin{problembox}[Q4-01\quad 计算图的构建与前向传播]
    \textbf{题目描述：}\\
    文档第 4.1 节介绍了计算图的概念。计算图是反向传播的"地图"——把数学表达式画成有向图，每个节点表示一次运算，每条边表示数据流动的方向。

    \vspace{0.3cm}
    考虑以下前向计算过程（一个简化的单样本损失计算）：
    \[
    z = x \cdot w + b, \quad \hat{y} = \sigma(z), \quad L = -\bigl(y \log \hat{y} + (1-y)\log(1-\hat{y})\bigr)
    \]
    其中 $x=2,\ w=0.5,\ b=1,\ y=1$（真实标签），$\sigma$ 为 Sigmoid 函数。

    \vspace{0.3cm}
    \textbf{问题：}
    \begin{enumerate}[(1)]
        \item 按照前向计算顺序，依次计算以下中间变量的数值（结果保留三位小数）：
        \begin{enumerate}[a.]
            \item $z = x \cdot w + b = $\underline{\hspace{4cm}}
            \item $\hat{y} = \sigma(z) = \dfrac{1}{1+e^{-z}} = $\underline{\hspace{3.5cm}}
            \item $L = -\bigl(y\log\hat{y} + (1-y)\log(1-\hat{y})\bigr) = $\underline{\hspace{3cm}}
        \end{enumerate}

        \item 计算图中每个节点只做\textbf{一种基本运算}。将上述计算分解，列出从输入 $(x, w, b, y)$ 到 $L$ 所经历的所有基本运算节点，并标注每个节点的输入和输出：

        \begin{center}
        \begin{tabular}{ccccc}
            \toprule
            节点编号 & 运算类型 & 输入变量 & 输出变量 & 输出数值 \\
            \midrule
            1 & 乘法 & $x,\ w$ & $t_1 = x \cdot w$ & \underline{\hspace{1.5cm}} \\
            2 & 加法 & $t_1,\ b$ & $t_2 = t_1 + b$ & \underline{\hspace{1.5cm}} \\
            3 & Sigmoid & $t_2$ & $\hat{y} = \sigma(t_2)$ & \underline{\hspace{1.5cm}} \\
            4 & 对数 & $\hat{y}$ & $t_3 = \log \hat{y}$ & \underline{\hspace{1.5cm}} \\
            5 & 乘法 & $y,\ t_3$ & $t_4 = y \cdot t_3$ & \underline{\hspace{1.5cm}} \\
            6 & 取负 & $t_4$ & $L = -t_4$ & \underline{\hspace{1.5cm}} \\
            \bottomrule
        \end{tabular}
        \end{center}
        （注：此处简化为 $y=1$ 时的情况，省略了 $(1-y)\log(1-\hat{y})$ 项）

        \item 计算图有两个核心优势：\textbf{模块化}和\textbf{局部计算}。用一句话解释"局部计算"是什么意思，以及它为什么让反向传播变得高效。
    \end{enumerate}

    \vspace{0.2cm}
    {\color{gray}\small\textit{来源溯源：[文档第 4 章 4.1 计算图]}}
\end{problembox}

\begin{beginnerbox}[小白必读：为什么要把计算"画"成图？]
    \begin{itemize}
        \item \textbf{计算图的本质}：把一个复杂的数学表达式拆解成一系列"只做一件事"的小节点，每个节点知道自己的输入是什么、输出是什么、以及这个运算的导数是多少。
        \item \textbf{为什么要拆？}拆成图之后，反向传播可以"逐节点"地传递梯度，每个节点只需要知道"上游传来的梯度"和"自己这个运算的局部导数"，乘一下就得到传给下游的梯度——不需要知道整个网络的全局结构。
        \item \textbf{计算图 = 自动微分的蓝图}：PyTorch 的 \texttt{autograd} 背后做的事，就是在前向传播时悄悄构建这张图，在 \texttt{backward()} 时沿图反向遍历，自动计算每个节点的梯度。
    \end{itemize}
\end{beginnerbox}

\begin{solutionbox}[参考答案与详细解析]
    \textbf{(1) 中间变量数值计算：}

    \begin{enumerate}[a.]
        \item $z = x \cdot w + b = 2 \times 0.5 + 1 = 1 + 1 = \boxed{2.000}$
        \item $\hat{y} = \sigma(2) = \dfrac{1}{1+e^{-2}} = \dfrac{1}{1+0.135} \approx \boxed{0.880}$
        \item $y=1$ 时，$(1-y)\log(1-\hat{y}) = 0$，故：
        \[
        L = -\log(0.880) \approx -(-0.128) = \boxed{0.128}
        \]
    \end{enumerate}

    \textbf{(2) 计算图节点数值：}

    \begin{center}
    \begin{tabular}{ccccc}
        \toprule
        节点编号 & 运算类型 & 输入变量 & 输出变量 & 输出数值 \\
        \midrule
        1 & 乘法 & $x=2,\ w=0.5$ & $t_1$ & $1.000$ \\
        2 & 加法 & $t_1=1,\ b=1$ & $t_2$ & $2.000$ \\
        3 & Sigmoid & $t_2=2$ & $\hat{y}$ & $0.880$ \\
        4 & 对数 & $\hat{y}=0.880$ & $t_3$ & $-0.128$ \\
        5 & 乘法 & $y=1,\ t_3=-0.128$ & $t_4$ & $-0.128$ \\
        6 & 取负 & $t_4=-0.128$ & $L$ & $0.128$ \\
        \bottomrule
    \end{tabular}
    \end{center}

    \textbf{(3) 局部计算的含义：}

    \textbf{局部计算}是指：计算图中每个节点在前向传播时只需要关注自己的输入，在反向传播时只需要知道上游传来的梯度和自己的局部导数（对自身输入的偏导数）——与网络其他部分的结构完全无关。这使得反向传播可以\textbf{分节点独立计算、模块化复用}，无论网络有多深多宽，每个节点的计算复杂度都是常数，整体复杂度仅与网络规模线性相关，远比直接对整个网络展开求导高效。
\end{solutionbox}


%% ============================================================
\Needspace{5\baselineskip}
\begin{problembox}[Q4-02\quad 链式法则：梯度沿计算图反向流动]
    \textbf{题目描述：}\\
    文档第 4.2 节介绍了链式法则（Chain Rule）——反向传播的数学基础。链式法则说的是：对于复合函数 $L = f(g(x))$，有：
    \[
    \frac{\partial L}{\partial x} = \frac{\partial L}{\partial g} \cdot \frac{\partial g}{\partial x}
    \]

    \vspace{0.3cm}
    考虑一个三层复合函数（对应一条三节点的计算图路径）：
    \[
    t_1 = 3x^2, \quad t_2 = t_1 + 5, \quad L = 2t_2
    \]

    \textbf{问题：}
    \begin{enumerate}[(1)]
        \item 用链式法则，推导 $\dfrac{\partial L}{\partial x}$。写出完整的推导过程，分三步展开：
        \[
        \frac{\partial L}{\partial x} = \frac{\partial L}{\partial t_2} \cdot \frac{\partial t_2}{\partial t_1} \cdot \frac{\partial t_1}{\partial x}
        \]

        \item 在计算图中，反向传播从 $L$ 开始，将梯度逐节点向左传递。设"上游梯度"为从右侧传来的梯度值，填写下表（取 $x=1$ 时的数值）：

        \begin{center}
        \begin{tabular}{p{1.5cm} p{3.5cm} p{3.5cm} p{3.0cm}}
            \toprule
            节点 & 该节点的局部导数 & 上游梯度（从右侧传来） & 传出的梯度（上游$\times$局部） \\
            \midrule
            $L=2t_2$ & $\dfrac{\partial L}{\partial t_2} = \underline{\quad}$ & $1$（Loss 对自身的导数） & $\underline{\hspace{2cm}}$ \\[1.5em]
            $t_2=t_1+5$ & $\dfrac{\partial t_2}{\partial t_1} = \underline{\quad}$ & \underline{\hspace{2cm}} & \underline{\hspace{2cm}} \\[1.5em]
            $t_1=3x^2$ & $\dfrac{\partial t_1}{\partial x} = \underline{\quad}$ & \underline{\hspace{2cm}} & \underline{\hspace{2cm}} \\
            \bottomrule
        \end{tabular}
        \end{center}

        \item 验证：直接计算 $L = 2(3x^2+5) = 6x^2+10$，对 $x$ 求导，在 $x=1$ 处的值是多少？与上表最终结果是否一致？

        \item 当计算图出现\textbf{分支}（即一个节点的输出同时被多个下游节点使用）时，该节点的梯度应该如何处理？（提示：考虑 $z = x + x^2$ 时，$x$ 出现了两次）
    \end{enumerate}

    \vspace{0.2cm}
    {\color{gray}\small\textit{来源溯源：[文档第 4 章 4.2 链式法则]}}
\end{problembox}

\begin{beginnerbox}[小白必读：链式法则为什么叫"链式"？]
    \begin{itemize}
        \item \textbf{链式}的含义：梯度像一条链条上的力，从最终输出端（Loss）出发，一环一环地向后传递，每经过一个节点就乘以该节点的局部导数，直到传递到每一个参数。
        \item \textbf{关键洞察}：每个节点不需要知道"Loss 从哪来"，只需要知道"上游告诉我梯度是多少"，然后用自己的局部导数乘一下，传给下游即可。这就是为什么可以把整个网络拆成独立模块——每个模块只实现自己的 \texttt{forward}（计算输出）和 \texttt{backward}（乘以局部导数传递梯度）。
        \item \textbf{链式法则的核心公式}：
        \[
        \frac{\partial L}{\partial x} = \underbrace{\frac{\partial L}{\partial y}}_{\text{上游梯度}} \times \underbrace{\frac{\partial y}{\partial x}}_{\text{局部导数}}
        \]
    \end{itemize}
\end{beginnerbox}

\begin{solutionbox}[参考答案与详细解析]
    \textbf{(1) 链式法则推导：}

    \[
    \frac{\partial L}{\partial t_2} = 2, \quad \frac{\partial t_2}{\partial t_1} = 1, \quad \frac{\partial t_1}{\partial x} = 6x
    \]
    \[
    \frac{\partial L}{\partial x} = 2 \times 1 \times 6x = 12x
    \]

    \textbf{(2) 反向传播逐节点填表（$x=1$）：}

    \begin{center}
    \begin{tabular}{p{1.5cm} p{3.0cm} p{3.5cm} p{3.0cm}}
        \toprule
        节点 & 局部导数 & 上游梯度 & 传出的梯度 \\
        \midrule
        $L=2t_2$ & $\partial L/\partial t_2 = 2$ & $1$ & $2 \times 1 = 2$ \\[0.8em]
        $t_2=t_1+5$ & $\partial t_2/\partial t_1 = 1$ & $2$ & $1 \times 2 = 2$ \\[0.8em]
        $t_1=3x^2$ & $\partial t_1/\partial x = 6x=6$ & $2$ & $6 \times 2 = \mathbf{12}$ \\
        \bottomrule
    \end{tabular}
    \end{center}

    \textbf{(3) 直接求导验证：}

    $L = 6x^2 + 10$，$\dfrac{\partial L}{\partial x} = 12x$，在 $x=1$ 处：$\dfrac{\partial L}{\partial x}\big|_{x=1} = 12$。

    与表中最终结果 $\mathbf{12}$ 完全一致。\checkmark

    \textbf{(4) 分支节点的梯度处理——梯度累加：}

    当一个节点的输出同时被多个下游节点使用时，该节点从多个下游接收梯度，应将所有来源的梯度\textbf{相加}。

    例如 $L = z = x + x^2$，可以理解为 $x$ 同时进入了加法节点（贡献 1）和平方节点（贡献 $2x$），两条路径的梯度相加：
    \[
    \frac{\partial L}{\partial x} = 1 + 2x
    \]
    在计算图中，这对应"两条从下游传来的梯度箭头在 $x$ 节点处汇合，执行求和"。

    \begin{keybox}[反向传播两大规则速记]
        \begin{itemize}
            \item \textbf{串联节点}（单路径）：梯度 = 上游梯度 $\times$ 局部导数（链式相乘）
            \item \textbf{分支节点}（多路径汇合）：梯度 = 所有下游传来梯度之\textbf{和}（梯度叠加）
        \end{itemize}
    \end{keybox}
\end{solutionbox}


%% ============================================================
\Needspace{5\baselineskip}
\begin{problembox}[Q4-03\quad 加法节点与乘法节点的反向传播]
    \textbf{题目描述：}\\
    文档第 4.3 节指出，任意复杂的计算图都可以分解为加法节点和乘法节点的组合。掌握这两种最基本节点的反向传播规则，是理解整个反向传播算法的基石。

    \vspace{0.3cm}
    \textbf{问题：}
    \begin{enumerate}[(1)]
        \item \textbf{加法节点}：设 $z = x + y$，上游传来的梯度为 $\dfrac{\partial L}{\partial z} = \delta$。

        推导 $\dfrac{\partial L}{\partial x}$ 和 $\dfrac{\partial L}{\partial y}$，并用一句话总结加法节点反向传播的规律。

        \item \textbf{乘法节点}：设 $z = x \cdot y$，上游传来的梯度为 $\dfrac{\partial L}{\partial z} = \delta$。

        推导 $\dfrac{\partial L}{\partial x}$ 和 $\dfrac{\partial L}{\partial y}$，并用一句话总结乘法节点反向传播的规律。

        \item \textbf{数值验证}：已知计算图中某乘法节点，前向传播时输入为 $x=3,\ y=-2$，上游梯度 $\delta = 4$。计算传给 $x$ 和 $y$ 的梯度。

        \item \textbf{综合练习}：对于表达式 $L = (x + y) \cdot z$，其中 $x=1,\ y=2,\ z=3$：
        \begin{enumerate}[a.]
            \item 写出前向传播的计算过程（设中间变量 $s = x + y$）；
            \item 从 $\dfrac{\partial L}{\partial L}=1$ 出发，逐节点反向传播，求出 $\dfrac{\partial L}{\partial x}$，$\dfrac{\partial L}{\partial y}$，$\dfrac{\partial L}{\partial z}$；
            \item 用直接求偏导数的方法验证答案（$L = xz + yz$，分别对 $x$、$y$、$z$ 求偏导）。
        \end{enumerate}
    \end{enumerate}

    \vspace{0.2cm}
    {\color{gray}\small\textit{来源溯源：[文档第 4 章 4.3 反向传播]}}
\end{problembox}

\begin{metaphorbox}[生活比喻：加法节点和乘法节点的梯度像什么？]
    \begin{itemize}
        \item \textbf{加法节点 = 快递分拣中心}：一个包裹到达，原封不动地分发给两个收件人（$x$ 和 $y$）。不管包裹大小，两边收到的都是同样大小的"副本"。梯度原样广播到所有输入。
        \item \textbf{乘法节点 = 权重放大器}：上游说"整体需要调 $\delta$ 倍"，但对 $x$ 的调整量由 $y$ 的大小决定（$\delta \cdot y$），对 $y$ 的调整量由 $x$ 的大小决定（$\delta \cdot x$）——谁对结果贡献越大（值越大），谁对另一方的梯度放大越多。
        \item \textbf{乘法节点的反直觉之处}：梯度中出现了\textbf{另一方}的值，这意味着如果某个输入值很大，对应的梯度也会很大——这正是为什么深度网络中权重初始化过大会引发梯度爆炸。
    \end{itemize}
\end{metaphorbox}

\begin{solutionbox}[参考答案与详细解析]
    \textbf{(1) 加法节点反向传播：}

    $z = x + y$，局部导数为 $\dfrac{\partial z}{\partial x} = 1$，$\dfrac{\partial z}{\partial y} = 1$。

    由链式法则：
    \[
    \frac{\partial L}{\partial x} = \frac{\partial L}{\partial z} \cdot \frac{\partial z}{\partial x} = \delta \cdot 1 = \delta
    \]
    \[
    \frac{\partial L}{\partial y} = \frac{\partial L}{\partial z} \cdot \frac{\partial z}{\partial y} = \delta \cdot 1 = \delta
    \]

    \textbf{规律}：加法节点将上游梯度\textbf{原样广播}给所有输入，不做任何放大或缩小。

    \textbf{(2) 乘法节点反向传播：}

    $z = x \cdot y$，局部导数为 $\dfrac{\partial z}{\partial x} = y$，$\dfrac{\partial z}{\partial y} = x$。

    \[
    \frac{\partial L}{\partial x} = \delta \cdot y, \quad \frac{\partial L}{\partial y} = \delta \cdot x
    \]

    \textbf{规律}：乘法节点将上游梯度乘以\textbf{另一方的前向传播值}后传出——$x$ 的梯度含 $y$，$y$ 的梯度含 $x$。

    \textbf{(3) 数值验证（$x=3,\ y=-2,\ \delta=4$）：}

    \[
    \frac{\partial L}{\partial x} = \delta \cdot y = 4 \times (-2) = \boxed{-8}
    \]
    \[
    \frac{\partial L}{\partial y} = \delta \cdot x = 4 \times 3 = \boxed{12}
    \]

    \textbf{(4) 综合练习（$L=(x+y)\cdot z$，$x=1,\ y=2,\ z=3$）：}

    \textbf{(a) 前向传播}：$s = x + y = 3$，$L = s \cdot z = 9$

    \textbf{(b) 反向传播}：

    从 $L = s \cdot z$ 出发，$\dfrac{\partial L}{\partial L}=1$：
    \[
    \frac{\partial L}{\partial s} = 1 \cdot z = 3, \quad \frac{\partial L}{\partial z} = 1 \cdot s = 3
    \]
    从 $s = x + y$ 出发，上游梯度为3：
    \[
    \frac{\partial L}{\partial x} = 3 \cdot 1 = \boxed{3}, \quad \frac{\partial L}{\partial y} = 3 \cdot 1 = \boxed{3}, \quad \frac{\partial L}{\partial z} = \boxed{3}
    \]

    \textbf{(c) 直接求偏导验证}：$L = xz + yz$

    \[
    \frac{\partial L}{\partial x} = z = 3\ \checkmark, \quad \frac{\partial L}{\partial y} = z = 3\ \checkmark, \quad \frac{\partial L}{\partial z} = x + y = 3\ \checkmark
    \]
    三者完全一致。
\end{solutionbox}


%% ============================================================
\Needspace{5\baselineskip}
\begin{problembox}[Q4-04\quad 激活层的反向传播：ReLU 与 Sigmoid]
    \textbf{题目描述：}\\
    文档第 4.4 节推导了 ReLU 和 Sigmoid 两种激活函数的反向传播公式。激活层是计算图中的"非线性节点"，其梯度的"干净程度"直接影响梯度能否顺利流向底层。

    \vspace{0.3cm}
    \textbf{问题：}
    \begin{enumerate}[(1)]
        \item \textbf{ReLU 的反向传播}：ReLU 定义为：
        \[
        y = \text{ReLU}(x) = \max(0, x)
        \]
        推导 $\dfrac{\partial L}{\partial x}$（用上游梯度 $\delta = \dfrac{\partial L}{\partial y}$ 表示），并解释为什么 ReLU 的反向传播可以用一个简单的"开关"来描述。

        \item \textbf{Sigmoid 的反向传播}：Sigmoid 定义为 $y = \sigma(x) = \dfrac{1}{1+e^{-x}}$，其导数有一个优美的形式：
        \[
        \frac{d\sigma}{dx} = \sigma(x)\bigl(1-\sigma(x)\bigr) = y(1-y)
        \]
        推导 $\dfrac{\partial L}{\partial x}$（用 $y$ 和 $\delta$ 表示），并计算：当 $y=0.880$（即上题 Q4-01 中的 $\hat{y}$），$\delta = 1$ 时，$\dfrac{\partial L}{\partial x}$ 的值（保留三位小数）。

        \item \textbf{梯度"干净程度"对比}：分析并填写下表：

        \begin{center}
        \begin{tabular}{p{2.0cm} p{4.0cm} p{4.5cm}}
            \toprule
            激活函数 & 反向传播公式 & 梯度特点（从防止梯度消失角度分析）\\
            \midrule
            ReLU & $\dfrac{\partial L}{\partial x} = \begin{cases} \delta & x > 0 \\ 0 & x \leq 0 \end{cases}$ & \underline{\hspace{4cm}} \\[2.5em]
            Sigmoid & $\dfrac{\partial L}{\partial x} = \delta \cdot y(1-y)$ & \underline{\hspace{4cm}} \\
            \bottomrule
        \end{tabular}
        \end{center}

        \item \textbf{代码实现}：用 Python/NumPy 实现 ReLU 层的 \texttt{forward} 和 \texttt{backward} 方法，要求 \texttt{backward} 利用 \texttt{forward} 阶段保存的信息：
        \begin{lstlisting}[language=Python]
import numpy as np

class ReLU:
    def __init__(self):
        self.mask = None   # 记录哪些位置被置零

    def forward(self, x):
        # 请补全
        pass

    def backward(self, dout):
        # 请补全（dout 是上游梯度）
        pass
        \end{lstlisting}
    \end{enumerate}

    \vspace{0.2cm}
    {\color{gray}\small\textit{来源溯源：[文档第 4 章 4.4 激活层的反向传播和实现]}}
\end{problembox}

\begin{beginnerbox}[小白必读：为什么激活函数的导数如此重要？]
    \begin{itemize}
        \item \textbf{激活层是梯度必经之路}：梯度从 Loss 出发，每经过一个激活层，就要乘以该激活函数的导数。如果导数很小（如 Sigmoid 在饱和区导数趋近于0），梯度经过多层后就会消失。
        \item \textbf{ReLU 的"开关"特性}：前向传播时，ReLU 把负输入变成0（开关关闭），正输入原样通过（开关打开）。反向传播时，开关关闭的神经元梯度直接变成0（那个方向的信号被彻底截断），开关打开的神经元梯度原样通过——干净利落，不引入额外衰减。
        \item \textbf{Sigmoid 导数的上限}：$y(1-y)$ 在 $y=0.5$ 时取最大值 $0.25$，且当 $y$ 接近0或1时（饱和区），导数趋近于0。这正是 Sigmoid 引发梯度消失的根本原因。
    \end{itemize}
\end{beginnerbox}

\begin{solutionbox}[参考答案与详细解析]
    \textbf{(1) ReLU 的反向传播：}

    前向传播：$y = \max(0, x)$，故局部导数为：
    \[
    \frac{\partial y}{\partial x} = \begin{cases} 1 & x > 0 \\ 0 & x \leq 0 \end{cases}
    \]
    由链式法则：
    \[
    \frac{\partial L}{\partial x} = \delta \cdot \frac{\partial y}{\partial x} = \begin{cases} \delta & x > 0 \\ 0 & x \leq 0 \end{cases}
    \]
    \textbf{"开关"描述}：前向时 $x>0$ 的位置"开关打开"，梯度原样通过；$x\leq0$ 的位置"开关关闭"，梯度直接置零。开关的状态（哪些位置为正）在前向传播时记录，反向传播时直接复用。

    \textbf{(2) Sigmoid 的反向传播：}

    \[
    \frac{\partial L}{\partial x} = \delta \cdot y(1-y)
    \]
    当 $y=0.880$，$\delta=1$ 时：
    \[
    \frac{\partial L}{\partial x} = 1 \times 0.880 \times (1-0.880) = 0.880 \times 0.120 = \boxed{0.106}
    \]

    \textbf{(3) 梯度特点对比：}

    \begin{center}
    \begin{tabular}{p{2.0cm} p{4.0cm} p{4.3cm}}
        \toprule
        激活函数 & 反向传播公式 & 梯度特点 \\
        \midrule
        ReLU & $\delta$（$x>0$）或 $0$（$x\leq0$） & 正值区域梯度恒为1，不引入额外衰减，可有效缓解梯度消失；但负值区域梯度恒为0（死亡ReLU风险） \\[1em]
        Sigmoid & $\delta \cdot y(1-y)$ & 导数最大值仅0.25，饱和区趋近于0，多层连乘后梯度快速消失，不适合深层网络隐藏层 \\
        \bottomrule
    \end{tabular}
    \end{center}

    \textbf{(4) ReLU 层代码实现：}

    \begin{lstlisting}[language=Python]
import numpy as np

class ReLU:
    def __init__(self):
        self.mask = None  # 记录前向传播中哪些位置 <= 0

    def forward(self, x):
        self.mask = (x <= 0)          # True 表示该位置被 ReLU 置零
        out = x.copy()
        out[self.mask] = 0            # 负值和零位置输出置零
        return out

    def backward(self, dout):
        dout[self.mask] = 0           # 前向被置零的位置，梯度也置零
        dx = dout
        return dx
    \end{lstlisting}

    \begin{keybox}[前向保存信息供反向使用——这是层实现的通用模式]
        \begin{itemize}
            \item \textbf{ReLU}：保存 \texttt{mask}（哪些位置 $\leq 0$），反向时直接用
            \item \textbf{Sigmoid}：保存前向输出 $y$，反向时计算 $y(1-y)$（无需重新算指数）
            \item \textbf{乘法节点}：分别保存 $x$ 和 $y$，反向时互换使用
        \end{itemize}
    \end{keybox}
\end{solutionbox}


%% ============================================================
\Needspace{5\baselineskip}
\begin{problembox}[Q4-05\quad Affine 层（全连接层）的反向传播]
    \textbf{题目描述：}\\
    文档第 4.5 节推导了 Affine 层（即全连接层）在\textbf{批量输入}下的反向传播。Affine 层是 $Y = XW + b$，当 $X$ 是一个批量矩阵时，梯度的推导需要特别注意形状。

    \vspace{0.3cm}
    设批量大小为 $N=2$，输入特征数 $D=3$，输出特征数 $M=4$：
    \[
    X_{(N \times D)} = \begin{pmatrix} 1 & 2 & 3 \\ 4 & 5 & 6 \end{pmatrix}, \quad
    W_{(D \times M)} = \begin{pmatrix} w_{11} & \cdots \\ \vdots & \ddots \end{pmatrix}, \quad
    b_{(1 \times M)} = \begin{pmatrix} b_1 & b_2 & b_3 & b_4 \end{pmatrix}
    \]
    前向传播：$Y = XW + b$，上游传来的梯度为 $\dfrac{\partial L}{\partial Y}$（形状与 $Y$ 相同，记为 $\delta_Y$，形状 $N \times M$）。

    \vspace{0.3cm}
    \textbf{问题：}
    \begin{enumerate}[(1)]
        \item 写出 $Y$ 的形状（$N \times M$）。

        \item 推导以下三个梯度（提示：维度分析是最快的推导方法——保证等式两边的形状一致）：
        \begin{enumerate}[a.]
            \item $\dfrac{\partial L}{\partial W}$（形状应为 $D \times M$，与 $W$ 相同）
            \item $\dfrac{\partial L}{\partial b}$（形状应为 $1 \times M$，与 $b$ 相同）
            \item $\dfrac{\partial L}{\partial X}$（形状应为 $N \times D$，与 $X$ 相同）
        \end{enumerate}

        \item \textbf{形状验证}：给出推导结果后，验证以下等式的形状一致性：
        \[
        \frac{\partial L}{\partial W} = X^T \delta_Y: \quad (D \times N)(N \times M) = D \times M\ \checkmark
        \]
        \[
        \frac{\partial L}{\partial X} = \delta_Y W^T: \quad (N \times M)(M \times D) = N \times D\ \checkmark
        \]
        \[
        \frac{\partial L}{\partial b} = \sum_{n=1}^{N}\delta_{Y,n}: \quad \text{对 }N\text{ 个样本的梯度求和，得 }1 \times M\ \checkmark
        \]

        \item \textbf{为什么 $b$ 的梯度要对 $N$ 个样本求和}？从偏置的广播机制和反向传播的梯度分支累加规则解释。

        \item 用 Python/NumPy 实现 Affine 层的 \texttt{forward} 和 \texttt{backward} 方法：
        \begin{lstlisting}[language=Python]
import numpy as np

class Affine:
    def __init__(self, W, b):
        self.W  = W
        self.b  = b
        self.x  = None   # 保存前向输入

    def forward(self, x):
        # 请补全
        pass

    def backward(self, dout):
        # 请补全，返回 dx，并更新 self.dW, self.db
        pass
        \end{lstlisting}
    \end{enumerate}

    \vspace{0.2cm}
    {\color{gray}\small\textit{来源溯源：[文档第 4 章 4.5 Affine 的反向传播和实现]}}
\end{problembox}

\begin{beginnerbox}[小白必读：矩阵求导怎么不出错？用形状来推！]
    \begin{itemize}
        \item \textbf{维度分析法}：对于矩阵运算的反向传播，最可靠的方法不是死记公式，而是\textbf{让等式两边的形状匹配}。例如，$\partial L / \partial W$ 的形状必须和 $W$ 一样是 $D \times M$，能凑出这个形状的只有 $X^T \delta_Y$（$(D \times N)(N \times M)$）。
        \item \textbf{转置的直觉}：矩阵乘法 $Y = XW$ 中，前向传播 $X$ 在左；反向传播计算 $\partial L/\partial W$ 时，$X$ 变成 $X^T$ 在左。这个"转置换位置"的模式在所有矩阵乘法的反向传播中都成立。
        \item \textbf{偏置的特殊处理}：偏置 $b$ 在前向传播时通过广播加到每个样本上（$N$ 次），相当于 $N$ 条路径共用同一个 $b$，反向传播时 $N$ 条路径的梯度需要全部加回给 $b$——所以要对样本维度求和。
    \end{itemize}
\end{beginnerbox}

\begin{solutionbox}[参考答案与详细解析]
    \textbf{(1) $Y$ 的形状：}

    $Y = XW + b$，$X$ 形状 $(N \times D) = (2 \times 3)$，$W$ 形状 $(D \times M) = (3 \times 4)$，故 $Y$ 形状为 $\boxed{N \times M = 2 \times 4}$。

    \textbf{(2) 三个梯度的推导：}

    \textbf{(a) $\partial L / \partial W$（形状 $D \times M$）：}

    由形状匹配，唯一能从 $X$（$(N \times D)$）和 $\delta_Y$（$(N \times M)$）凑出 $D \times M$ 的是：
    \[
    \frac{\partial L}{\partial W} = X^T \delta_Y \quad (D \times N)(N \times M) = D \times M
    \]

    \textbf{(b) $\partial L / \partial b$（形状 $1 \times M$）：}

    $b$ 对每个样本都广播了一次，每次前向传播产生一条梯度路径，反向时求和：
    \[
    \frac{\partial L}{\partial b} = \sum_{n=1}^{N} \delta_{Y,n} = \mathbf{1}^T \delta_Y \quad (1 \times N)(N \times M) = 1 \times M
    \]
    即对 $\delta_Y$ 按行求和（\texttt{np.sum(dout, axis=0)}）。

    \textbf{(c) $\partial L / \partial X$（形状 $N \times D$）：}

    \[
    \frac{\partial L}{\partial X} = \delta_Y W^T \quad (N \times M)(M \times D) = N \times D
    \]

    \textbf{(3) 形状验证}：已在题目中给出，三式均验证通过。\checkmark

    \textbf{(4) 偏置梯度求和的原因：}

    前向传播中，$b$ 通过广播同时加到了 $N$ 个样本的输出上，相当于在计算图中，$b$ 这个节点同时连接到了 $N$ 个"加法节点"（每个样本对应一个）。根据分支节点的梯度累加规则，$b$ 从 $N$ 条路径各接收一份梯度，总梯度 = 各路径梯度之\textbf{和}，即对批量维度（axis=0）求和。

    \textbf{(5) Affine 层代码实现：}

    \begin{lstlisting}[language=Python]
import numpy as np

class Affine:
    def __init__(self, W, b):
        self.W  = W
        self.b  = b
        self.x  = None
        self.dW = None
        self.db = None

    def forward(self, x):
        self.x = x                          # 保存输入供反向使用
        out = x @ self.W + self.b           # Y = XW + b
        return out

    def backward(self, dout):
        dx       = dout @ self.W.T          # dL/dX = δ W^T
        self.dW  = self.x.T @ dout          # dL/dW = X^T δ
        self.db  = np.sum(dout, axis=0)     # dL/db = 各样本梯度求和
        return dx
    \end{lstlisting}
\end{solutionbox}


%% ============================================================
\Needspace{5\baselineskip}
\begin{problembox}[Q4-06\quad 综合代码实现：两层网络的完整反向传播]
    \textbf{题目描述：}\\
    文档第 4.7 节给出了一个完整的反向传播综合实现。以下是一个简化的两层神经网络结构，包含 Affine 层、ReLU 激活层和 Softmax-CrossEntropy 输出层：

    \begin{lstlisting}[language=Python]
class TwoLayerNet:
    def __init__(self, input_size, hidden_size, output_size):
        # 初始化权重（已省略）
        self.layers = {
            'Affine1': Affine(W1, b1),
            'ReLU1':   ReLU(),
            'Affine2': Affine(W2, b2),
        }
        self.last_layer = SoftmaxWithLoss()

    def predict(self, x):
        for layer in self.layers.values():
            x = layer.forward(x)
        return x

    def loss(self, x, t):
        y = self.predict(x)
        return self.last_layer.forward(y, t)

    def gradient(self, x, t):
        # 前向传播
        self.loss(x, t)

        # 反向传播
        dout = 1                                    # 步骤A
        dout = self.last_layer.backward(dout)       # 步骤B
        layers = list(self.layers.values())
        layers.reverse()
        for layer in layers:
            dout = layer.backward(dout)             # 步骤C

        # 提取各层梯度
        grads = {
            'W1': self.layers['Affine1'].dW,
            'b1': self.layers['Affine1'].db,
            'W2': self.layers['Affine2'].dW,
            'b2': self.layers['Affine2'].db,
        }
        return grads
    \end{lstlisting}

    \textbf{问题：}
    \begin{enumerate}[(1)]
        \item \textbf{前向传播数据流}：\texttt{predict} 方法中，数据依次经过哪些层？从输入 $x$ 到最终 logits 的变换链是什么（用符号表示）？

        \item \textbf{步骤A 的含义}：\texttt{dout = 1} 代表什么含义？为什么从 $\dfrac{\partial L}{\partial L} = 1$ 开始反向传播？

        \item \textbf{步骤B 的 SoftmaxWithLoss 层}：该层将 Softmax 和 CrossEntropy 合并。其反向传播有一个极其简洁的结论：
        \[
        \frac{\partial L}{\partial z_k} = y_k - t_k
        \]
        其中 $y_k$ 是 Softmax 输出概率，$t_k$ 是真实标签（one-hot）。解释这个公式的直觉含义：当预测概率 $y_k$ 和真实标签 $t_k$ 完全吻合时，梯度是多少？这意味着什么？

        \item \textbf{步骤C 的反向顺序}：\texttt{layers.reverse()} 将层列表反转后再依次调用 \texttt{backward}。解释为什么反向传播必须按照与前向传播相\textbf{反}的顺序执行。

        \item \textbf{梯度数值验证}：以下代码用"数值微分"法（有限差分近似）验证反向传播计算的梯度是否正确，请解释其原理，并说明为什么实际训练中不用数值微分而用反向传播：
        \begin{lstlisting}[language=Python]
def numerical_gradient(f, x):
    h = 1e-4
    grad = np.zeros_like(x)
    for i in range(x.size):
        x_flat = x.flatten()
        x_flat[i] += h
        fxh1 = f(x_flat.reshape(x.shape))
        x_flat[i] -= 2 * h
        fxh2 = f(x_flat.reshape(x.shape))
        grad.flat[i] = (fxh1 - fxh2) / (2 * h)
        x_flat[i] += h
    return grad
        \end{lstlisting}
    \end{enumerate}

    \vspace{0.2cm}
    {\color{gray}\small\textit{来源溯源：[文档第 4 章 4.6 输出层的反向传播和实现 / 4.7 反向传播的综合代码实现]}}
\end{problembox}

\begin{beginnerbox}[小白必读：为什么把 Softmax 和 CrossEntropy 合并成一个层？]
    \begin{itemize}
        \item \textbf{数值稳定性}：Softmax 分子分母中有 $e^{z_k}$，当 $z_k$ 很大时会数值溢出。合并后可以用"log-sum-exp"技巧在对数域计算，避免中间步骤的数值溢出。
        \item \textbf{梯度极简}：分开计算时，CrossEntropy 对 Softmax 输出的梯度要乘以 Softmax 自身的雅可比矩阵，形式复杂。合并后，两者的复杂项恰好约分，最终梯度简化为 $y_k - t_k$，既简洁又高效。
        \item \textbf{实践中对应}：PyTorch 的 \texttt{nn.CrossEntropyLoss} 就是 \texttt{LogSoftmax + NLLLoss} 的合并版，必须接收 logits（未经 Softmax 的原始输出）而不是概率，原因正在于此。
    \end{itemize}
\end{beginnerbox}

\begin{solutionbox}[参考答案与详细解析]
    \textbf{(1) 前向传播数据流：}

    数据依次经过：\texttt{Affine1} $\to$ \texttt{ReLU1} $\to$ \texttt{Affine2}。符号表示：
    \[
    x \xrightarrow{XW_1+b_1} a_1 \xrightarrow{\text{ReLU}} h_1 \xrightarrow{h_1W_2+b_2} \text{logits}
    \]
    \texttt{loss} 方法额外经过 \texttt{SoftmaxWithLoss}，计算最终损失值。

    \textbf{(2) 步骤A \texttt{dout=1} 的含义：}

    $\dfrac{\partial L}{\partial L} = 1$ 是恒等式——Loss 对自身的导数等于1。这是反向传播的起点：最终标量 Loss 对自身的梯度为1，相当于告诉第一个反向节点"整体的权重是1"，之后链式法则会自动将这个"1"按因果比例分配给每一层。若将 Loss 乘以某个系数（如取平均），这里就应该填入相应系数。

    \textbf{(3) SoftmaxWithLoss 反向传播的直觉含义：}

    梯度 $y_k - t_k$ 的直觉：
    \begin{itemize}
        \item 若 $y_k = t_k$（预测完全正确）：梯度 $= 0$，无需调整；
        \item 若 $y_k > t_k$（预测概率偏高）：梯度 $> 0$，应减小 $z_k$（降低该类别的"分数"）；
        \item 若 $y_k < t_k$（预测概率偏低）：梯度 $< 0$，应增大 $z_k$（提高该类别的"分数"）。
    \end{itemize}
    当预测与标签完全吻合时，梯度为\textbf{零向量}，参数不再更新——模型已经找到了"完美答案"，停止学习，这正是我们期望的行为。

    \textbf{(4) 反向顺序的必要性：}

    反向传播是链式法则的逐节点展开。计算某层的梯度时，需要"上游（更接近 Loss 的方向）已经算好的梯度"。若不按反向顺序，某层的上游梯度尚未计算完毕，就无法得到正确的传入梯度，链式法则断裂。因此，必须从 Loss 出发，严格按照前向传播的逆序依次调用 \texttt{backward}。

    \textbf{(5) 数值微分的原理及局限：}

    数值微分利用导数定义的近似（中心差分，误差 $O(h^2)$）：
    \[
    \frac{\partial f}{\partial x_i} \approx \frac{f(\ldots, x_i+h, \ldots) - f(\ldots, x_i-h, \ldots)}{2h}
    \]
    代码对每个参数 $x_i$ 分别加减一个小量 $h=10^{-4}$，调用两次前向传播，用差商近似偏导数。

    \textbf{为什么实际训练不用数值微分}：
    \begin{itemize}
        \item \textbf{计算量极大}：对每个参数都需要两次前向传播。一个拥有 $10^7$ 个参数的网络，需要 $2 \times 10^7$ 次前向传播——实际训练中不可接受；
        \item \textbf{反向传播只需一次}：一次前向 + 一次反向传播，就能同时算出所有参数的梯度，效率高出数量级；
        \item \textbf{数值微分的正确用途}：\textbf{梯度检验（Gradient Check）}——在开发新层时，用数值微分验证反向传播实现是否正确，而非用于实际训练。
    \end{itemize}

    \begin{keybox}[梯度检验的工作流程]
        \begin{enumerate}[(1)]
            \item 实现层的 \texttt{forward} 和 \texttt{backward}；
            \item 用数值微分计算"近似梯度"；
            \item 用 \texttt{backward} 计算"解析梯度"；
            \item 比较两者差异（通常要求相对误差 $< 10^{-5}$）；
            \item 若相差很大，说明 \texttt{backward} 实现有 bug。
        \end{enumerate}
    \end{keybox}
\end{solutionbox}

\begin{appbox}[现实应用场景]
    \textbf{反向传播——60年前的思想，今天的核心引擎：}\\
    反向传播算法的核心思想最早可追溯到1970年代，但直到1986年被引入神经网络后才真正改变世界。今天，PyTorch、TensorFlow 等框架的自动微分系统，本质上都是"动态计算图 + 反向传播"——你写的每一行前向计算代码，框架都在后台悄悄记录计算图，一旦调用 \texttt{.backward()}，梯度就沿图反向流动，自动填入每个参数的 \texttt{.grad}。理解本章，相当于理解了深度学习框架最核心的那台"发动机"是如何工作的。
\end{appbox}

% ===== 本章小结 =====
\section{本章小结：反向传播的全景}

学完这 6 道题，如果还觉得"原理飘在空中"，试试把整章内容浓缩成这几句画面：

\begin{itemize}
    \item \textbf{计算图} = 把复杂计算拆成只做一件事的小节点，每个节点只需要知道自己的局部导数；
    \item \textbf{链式法则} = 梯度的传递方式：上游梯度 $\times$ 局部导数 = 传给下游的梯度；分支处梯度相加；
    \item \textbf{加法节点} = 梯度原样广播；\textbf{乘法节点} = 梯度乘以对方的前向值（互换）；
    \item \textbf{ReLU} = 梯度的开关，正值区原样通过，负值区直接截断；\textbf{Sigmoid} = 梯度最大只有 $0.25$，多层后迅速消失；
    \item \textbf{Affine 层} = $\partial L/\partial W = X^T\delta_Y$，$\partial L/\partial X = \delta_Y W^T$，$\partial L/\partial b$ = 对样本维度求和；
    \item \textbf{完整流程} = 前向算 Loss，反向从1出发按逆序逐层调用 \texttt{backward}，梯度自动流到每个参数。
\end{itemize}

\begin{metaphorbox}[给零基础读者的一句总比喻]
    \textbf{如果把神经网络比作一条生产流水线：}
    \begin{itemize}
        \item \textbf{前向传播} = 原材料从左到右经过每道工序，最终生产出产品；
        \item \textbf{Loss} = 质检员打出的"不合格分数"；
        \item \textbf{反向传播} = "责任追溯"——质检报告从右到左传递，每道工序根据自己的"参与比例"（局部导数），承担相应的责任（梯度）；
        \item \textbf{梯度下降} = 每道工序根据自己承担的责任，调整操作参数，下次少出错；
        \item \textbf{自动微分框架} = 一套自动记录每道工序参与比例的管理系统——你只需要设计工序（前向计算），责任追溯（反向传播）全自动完成。
    \end{itemize}
\end{metaphorbox}

\end{document}
